\begingroup
\section*{Material complementario}
\addcontentsline{toc}{section}{Material complementario}
\begin{apareferences}
\aparef{ref:comparacion-bases}{Centro Javeriano de Competitividad. (2026, 16 de noviembre). \textit{Comparación de las propuestas del CJC con las Bases del PND y mecanismos de ejecución}. GitHub. \url{https://github.com/oliverpardo1979/Informe-CJC-PND/blob/main/material_complementario/Comparacion_Bases_PND_2026-11-13.md}}
\end{apareferences}

El cuadro identifica los compromisos de las Bases, los aportes del CJC, las vías de ejecución y las decisiones que cambiarían.

\begin{apareferences}
\aparef{ref:matriz-jorge}{Orozco, J. (s. f.). \textit{Matriz de políticas e iniciativas de inteligencia artificial} [Matriz de datos, versión consultada el 21 de octubre de 2026]. Centro Javeriano de Competitividad. \url{https://github.com/oliverpardo1979/Informe-CJC-PND/blob/main/material_complementario/Matriz_politicas_IA_Jorge_Orozco_2026-10-21.xlsm}}
\end{apareferences}

La matriz organiza experiencias internacionales por país y actividad económica, describe los instrumentos y su alcance de implementación y enlaza las fuentes oficiales. Incluye las hojas Matriz evidencia, Resumen y Criterios.

\begin{apareferences}
\aparef{ref:informe-ia}{Pardo, O., \& Orozco, J. (2026). \textit{Inteligencia artificial y empleo en Colombia} (Informe 03). Centro Javeriano de Competitividad, Pontificia Universidad Javeriana. \url{https://github.com/oliverpardo1979/informe_ia}}
\end{apareferences}

El repositorio contiene el informe, su metodología y los datos.

\textbf{Otros documentos de consulta.}
\begin{apareferences}
\aparef{ref:bib:mincit_fabricas_2023}{Ministerio de Comercio, Industria y Turismo. (2023, 20 de junio). \textit{En el Día Mundial de la Productividad, Gobierno nacional presenta programa para que 2.250 empresas sean más productivas y sostenibles}. \url{https://www.mincit.gov.co/prensa/noticias/industria/fabricas-de-productividad-y-sostenibilidad-2023}}

\end{apareferences}
\par\endgroup
